\documentclass[11pt,compress,t,notes=noshow, xcolor=table]{beamer}

\input{../../style/preamble}
\input{../../latex-math/basic-math}
\input{../../latex-math/basic-ml}
\input{../../latex-math/optim-basics}

\title{Optimization in Machine Learning}

\begin{document}

\titlemeta{
Constrained Optimization
}{
Penalty-Based and Primal-Dual Methods
}{
figure/primal_dual_penalty_barrier.pdf
}{
\item Quadratic penalties for violated constraints
\item Log-barrier and interior-point ideas
\item Why penalty / barrier parameters must change over time
\item Perturbed KKT systems and primal-dual interior-point updates
}

\begin{framev}[fs=scriptsize]{From Duality to Hybrid Methods}
\textbf{Starting point:}
\[
\min_{\xv \in \R^d} F(\xv)
\qquad
\text{s.t. } \qquad
g_i(\xv) \le 0,\;
h_j(\xv)=0.
\]
\begin{itemize}
\item Projection kept iterates feasible; duality relaxed constraints through multipliers.
\item \textbf{Penalty methods} allow violations but charge them inside the objective.
\item \textbf{Barrier methods} stay strictly feasible and approach the boundary from the interior.
\end{itemize}
\imageC[1]{figure/primal_dual_penalty_barrier.pdf}
\end{framev}


\begin{framev}[fs=scriptsize]{Penalty Path on the Box Toy}
Reuse
\[
\min_{\wv}
F(\wv)
=
(w_1 - 1.2)^2
+
(w_1 - 1.2)(w_2 - 0.2)
+
(w_2 - 0.2)^2
\]
\[
\text{s.t. }
\qquad
0 \le w_1 \le 1,
\qquad
0 \le w_2 \le 0.9.
\]
\splitV[0.58]{
Quadratic penalty relaxation:
\[
R(\wv,\rho)
:=
F(\wv)
+
\frac{\rho}{2}\sum_{i=1}^4 \max\{g_i(\wv),0\}^2.
\]
\begin{itemize}
\item The unconstrained minimizer $(1.2,0.2)$ violates only the upper bound $w_1 \le 1$.
\item For moderate $\rho$, the relaxed minimizer stays slightly outside the box.
\item As $\rho$ grows, the penalty path approaches the constrained optimum $(1,0.3)$ from the infeasible side.
\end{itemize}
}{
\image[1]{figure/primal_dual_penalty_path.pdf}
}
\end{framev}


\begin{framei}[fs=small]{Exact Versus Approximate Inner Solves}
\item In penalty methods, feasible points receive \textbf{no extra cost} beyond the original objective.
\item Exact statement: if the exact minimizer of the relaxed problem is feasible for the original constraints, then it also solves the original constrained problem.
\item In practice, each relaxed problem is solved only approximately.
\item Therefore, feasibility of the returned inner iterate by itself is \textbf{not} a proof of optimality for the original constrained problem.
\item A safe continuation strategy is:
\begin{enumerate}
\item choose an initial penalty weight $\rho_0 > 0$
\item minimize $R(\wv,\rho_t)$ to sufficient inner accuracy
\item if the returned point is feasible and the inner solve is accurate enough, we may stop
\item otherwise increase $\rho_t$ and solve again
\end{enumerate}
\item A common choice is geometric growth, for example $\rho_{t+1} = 2 \rho_t$.
\item The solution from the previous outer step is a natural warm start for the next relaxed problem.
\end{framei}


\begin{framei}[fs=footnotesize]{General Quadratic Penalty Relaxation}
\item Consider the constrained problem
$$
\min_{\xv} F(\xv)
\qquad
\text{s.t. }
\qquad
g_i(\xv) \le 0,\; i=1,\dots,m,
\qquad
h_j(\xv) = 0,\; j=1,\dots,k.
$$
\item A standard quadratic penalty relaxation is
$$
R(\xv,\rho)
:=
F(\xv)
+
\frac{\rho}{2}
\left(
\sum_{i=1}^m \max\{0,g_i(\xv)\}^2
+
\sum_{j=1}^k h_j(\xv)^2
\right),
\qquad \rho > 0.
$$
\item Inequality constraints are penalized only when violated; equality constraints are penalized whenever they are nonzero.
\item Its gradient is
$$
\nabla_{\xv} R(\xv,\rho)
=
\nabla F(\xv)
+
\rho \sum_{i=1}^m \max\{g_i(\xv),0\}\nabla g_i(\xv)
+
\rho \sum_{j=1}^k h_j(\xv)\nabla h_j(\xv).
$$
\end{framei}


\begin{framei}[fs=footnotesize]{Penalty Methods: Practical Behavior}
\item Penalty relaxations are attractive because they let us reuse unconstrained solvers such as GD, coordinate descent, or Newton-type methods.
\item However, the iterates of the relaxed problem need \textbf{not} be feasible while $\rho$ is still moderate.
\item Choosing an enormous penalty weight from the start is usually a bad idea:
\begin{itemize}
\item the penalty term dominates the geometry
\item the relaxed objective becomes ill-conditioned
\item descent methods can bounce or converge very slowly
\end{itemize}
\item Therefore, one usually starts with a smaller $\rho$, solves a reasonably conditioned relaxed problem, and then increases $\rho$ gradually.
\item Exact feasibility is often only approached in later outer iterations, once the penalty weight is large enough.
\end{framei}


\begin{framev}[fs=scriptsize]{Barrier Path on the Box Toy}
Barrier methods keep the inequality constraints explicit and replace the boundary by an infinite wall.
\begin{itemize}
\item For the same box toy,
\[
B(\wv,\mu)
:=
F(\wv)
- \mu \Big(
\log w_1 + \log(1-w_1) + \log w_2 + \log(0.9-w_2)
\Big).
\]
\item It is defined only for \textbf{strictly feasible} points:
\[
0 < w_1 < 1,
\qquad
0 < w_2 < 0.9.
\]
\item As $\mu \downarrow 0$, the barrier path approaches $(1,0.3)$ from the interior.
\end{itemize}
\imageC[0.5]{figure/primal_dual_barrier_path.pdf}
\end{framev}


\begin{framei}[fs=small]{Barrier Gradient and Continuation}
\item For the logarithmic barrier, the gradient is
$$
\nabla_{\xv} B(\xv,\mu)
=
\nabla F(\xv)
+
\mu \sum_{i=1}^m \frac{\nabla g_i(\xv)}{-g_i(\xv)}.
$$
\item As soon as one approaches the boundary $g_i(\xv) \to 0^-$, the term $\nabla g_i(\xv) / (-g_i(\xv))$ becomes very large.
\item This is the barrier effect: the objective rises sharply near the boundary and prevents infeasible steps.
\item The continuation strategy is the reverse of penalty methods:
\begin{itemize}
\item start with a relatively large barrier weight $\mu$
\item solve the barrier problem
\item reduce $\mu$ and warm-start the next solve
\end{itemize}
\item Large $\mu$ keeps the iterate well inside the feasible region; small $\mu$ allows it to move closer to the true constrained optimum.
\end{framei}


\begin{framev}[fs=scriptsize]{Barrier Stationarity and the Central Path}
With slacks $s_i := -g_i(\xv) > 0$ and multipliers $\lambda_i := \mu / s_i > 0$,
barrier stationarity becomes
$$
\nabla F(\xv) + \sum_{i=1}^m \lambda_i \nabla g_i(\xv) = \zero,
\qquad
\lambda_i s_i = \mu.
$$
\begin{itemize}
\item This is complementary slackness perturbed away from zero.
\item The \textbf{central path} keeps slack and multiplier both positive while $\mu > 0$.
\end{itemize}
\imageC[0.65]{figure/primal_dual_central_path.pdf}
\end{framev}


\begin{framei}[fs=footnotesize]{Primal-Dual Interior-Point View}
\item Interior-point methods solve the primal and dual variables \textbf{together} through the perturbed KKT system
$$
\nabla F(\xv)
+
\sum_{i=1}^m \lambda_i \nabla g_i(\xv)
+
\sum_{j=1}^k \nu_j \nabla h_j(\xv)
=
\zero,
$$
$$
g_i(\xv) + s_i = 0,
\qquad
h_j(\xv)=0,
\qquad
s_i > 0,\;
\lambda_i > 0,
$$
$$
\lambda_i s_i = \mu,
\qquad i=1,\dots,m.
$$
\item Newton steps are used to solve these equations approximately for a fixed $\mu$.
\item Then $\mu$ is decreased and the next solve is warm-started from the current primal-dual pair.
\item This is the clean primal-dual picture: the multipliers, slacks, and primal variables are coupled by one system instead of by an ad-hoc update rule.
\end{framei}


\begin{framev}{Take-Away}
\begin{itemize}
\item \textbf{Penalty methods:} add violation costs and increase $\rho$ until feasibility is reached or closely approximated.
\item \textbf{Barrier methods:} keep iterates strictly feasible and decrease $\mu$ to approach the optimum from the interior.
\item \textbf{Primal-dual view:} barrier stationarity already contains multiplier information and leads directly to perturbed KKT conditions.
\item \textbf{Trade-off:} penalties are simple but allow infeasible iterates; barriers preserve feasibility but require an interior start and careful continuation.
\end{itemize}
\end{framev}

\endlecture
\end{document}
